%% Individual Task 1: Part 1 -- Case Studies in Data Science (COSC2669/COSC2816)
%% Based on the ACM "acmart" class, manuscript format.

\documentclass[manuscript,screen,review,nonacm]{acmart}

\settopmatter{printacmref=false} % Removes citation information below abstract
\renewcommand\footnotetextcopyrightpermission[1]{} % removes footnote with conference information in first column
\pagestyle{plain} % removes running headers

%% Fonts used in the template cannot be substituted; margin
%% adjustments are not allowed.
%%
%% \BibTeX command to typeset BibTeX logo in the docs
\AtBeginDocument{%
  \providecommand\BibTeX{{%
    \normalfont B\kern-0.5em{\scshape i\kern-0.25em b}\kern-0.8em\TeX}}}

%Add any additional packages and commands here

\newcommand{\wordcount}[1]{\noindent\textbf{[Word count: #1]}} %macro for word count

\newcommand{\todo}[1]{\textcolor{red}{\bf [[TODO: #1]]}}
\newcommand{\instructions}[1]{\textcolor{blue}{\textbf{[[Remove before submitting:} #1]]}} %Instructions to be removed before submitting

\newcommand{\update}[1]{\textcolor{blue}{\textbf{[[Update before submitting:} #1]]}} %Instructions to be removed before submitting

%\newcommand{\instructions}[1]{} %Instructions to be removed before submitting

%extra packages for the tables in the appendix
\usepackage{longtable}
\usepackage{array}
\renewcommand{\arraystretch}{1.5}

%--- additions for Individual Task 1 -------------------------------------
\usepackage{listings}          % source-code listings in the appendix
\usepackage{multirow}          % merged cells in the results tables

\lstdefinestyle{pycode}{
  language=Python,
  basicstyle=\ttfamily\scriptsize,
  keywordstyle=\color{ACMPurple},
  commentstyle=\color{ACMGreen}\itshape,
  stringstyle=\color{ACMDarkBlue},
  numbers=left,
  numberstyle=\tiny\color{gray},
  numbersep=6pt,
  frame=single,
  rulecolor=\color{gray!40},
  breaklines=true,
  breakatwhitespace=true,
  showstringspaces=false,
  columns=fullflexible,
  keepspaces=true,
  tabsize=4,
  captionpos=b,
  xleftmargin=12pt,
}
\lstset{style=pycode}


\begin{document}

\title{Individual Task 1: Part 1}

\author{Heet Chanchad (s4218211)}
\affiliation{%
  \institution{RMIT University}
  \city{Melbourne}
  \country{Australia}}
\email{s4218211@student.rmit.edu.au}

\renewcommand{\shortauthors}{Chanchad}

\begin{abstract}
  This report examines the role of Machine Learning Engineer in Trust and Safety
  Engineering, using a live graduate advertisement at TikTok as its anchor. The
  professional problem the role exists to solve is triage under a fixed review budget: a
  platform receives far more user-generated content each day than any moderation team can
  read, so a classifier must rank items by the probability of violation and spend scarce
  human attention where it changes an outcome. Two public corpora of crowd-labelled user
  text are used to study that problem --- 24,783 annotated tweets and 159,686 annotated
  Wikipedia Talk comments --- and two machine learning algorithms, a linear support vector
  machine and a feed-forward neural network, are compared under stratified cross
  validation. Because the removal-worthy class covers only 5.8\% and 9.6\% of the two
  corpora respectively, the evaluation deliberately avoids accuracy in favour of
  precision, recall, area under the precision-recall curve, and precision over the slice
  of the ranked stream a fixed review team could actually process. The analysis finds that
  the two data sources are complementary in what they teach about model choice but
  directly contradictory in what they teach about labels, and that this contradiction is
  invisible to any single-dataset evaluation.

  % the # must be escaped here: the abstract body is stored in a macro first,
  % where a bare # is read as a parameter token and breaks the build.
  \textbf{Overleaf project URL:} \url{https://www.overleaf.com/read/swdgzkxnjjst\#264b95}
\end{abstract}

\maketitle

%----------------------------------
\section{Data Science Role}
\label{sec:Data Science Role}
%----------------------------------

%%----------------------------------
\subsection{Job Description and Role Identification}
\label{subsec:Job Description and Role Identification}
%----------------------------------

The role selected is \textbf{Machine Learning Engineer Graduate (Trust and Safety
Engineering) -- 2027 Start} at \textbf{TikTok}, based in Sydney, New South Wales, and
open to international applicants. Applications close 30 August 2026. The full
advertisement is reproduced in Appendix~\ref{sec:Job Listing}.

The advertisement places the engineer inside TikTok's Trust \& Safety team, which
describes itself as having ``invested heavily in human and machine-based moderation to
remove harmful content quickly and often before it reaches our general community''. The
listed responsibilities are explicitly build-oriented: to ``develop and build up
highly-scalable classifiers, tools, models and algorithms leveraging cutting-edge machine
learning, computer vision and data mining technologies'', to ``improve our trust and
safety strategy and work on model iterations'', and to ``develop and implement innovative
machine learning algorithms to manage business risks''. Minimum qualifications are a
Master's or PhD in computer science or a related engineering field, solid knowledge of at
least one of machine learning, pattern recognition, natural language processing, data
mining or computer vision, and a firm understanding of data structures and algorithms.

The professional problem behind the role is one of \emph{scale under a fixed review
budget}. A platform of TikTok's size receives far more user-generated content each day
than any moderation team can read, yet removing content automatically and without recourse
is unacceptable to creators and regulators alike. Neither exhaustive human review nor
fully autonomous removal is viable. The engineering problem is therefore triage: rank
incoming items by the probability that they violate policy, so that a finite number of
human moderators spend their attention where it changes an outcome. This is the problem
the remainder of this report addresses, and it is what makes the role a machine learning
\emph{engineering} role rather than an analytics one --- the deliverable is a ranked queue
running in production, monitored and re-trained, not a slide deck.

\textbf{Job listing URL:} \url{https://au.gradconnection.com/employers/tiktok/jobs/tiktok-machine-learning-engineer-graduate-trust-and-safety-engineering-2027-start-5/}

%%----------------------------------
\subsection{Industry Context}
\label{subsec:Industry Context}
%----------------------------------

Short-form video platforms cannot human-review millions of daily uploads. The trust and
safety field uses data science to detect hate speech, harassment, spam and coordinated
abuse, to rank what scarce moderators see first, and to monitor whether enforcement stays
consistent across languages, dialects and regions.

\wordcount{45}

%%----------------------------------
\subsection{Role Focus}
\label{subsec:Role Focus}
%----------------------------------

Prediction, delivered as production software. The advertisement asks the engineer to
``develop and build up highly-scalable classifiers, tools, models and algorithms'' and to
``work on model iterations'' to ``manage business risks''. Insight generation is
secondary; the deliverable is a deployed classifier that ranks content for moderator
review.

\wordcount{47}

%%----------------------------------
\subsection{Value Creation}
\label{subsec:Value Creation}
%----------------------------------

Automated triage lets a fixed moderation team cover an expanding content stream, so
harmful material is removed before it spreads. Better ranking also protects creators from
wrongful removal and advertisers from unsafe placement, while producing the enforcement
evidence regulators such as Australia's eSafety Commissioner require~\cite{esafety2022bosa}.

\wordcount{46}

%%----------------------------------
\subsection{Values and Culture Alignment}
\label{subsec:Values and Culture Alignment}
%----------------------------------

The team frames its mission as keeping people ``safe and empowered to create'', which
matches how I want to use machine learning: reducing harm rather than maximising
engagement. I value measured claims over confident ones, and this role's emphasis on
experiments, validation and iteration rewards that habit.

\wordcount{47}

%%----------------------------------
\subsection{Diversity and Inclusion}
\label{subsec:Diversity and Inclusion}
%----------------------------------

Moderation models inherit their annotators' judgements. The Wikipedia corpus used here
ships annotator demographics precisely because who labels the data shapes what counts as
toxic~\cite{wulczyn2017ex}. Diverse engineering teams are likelier to disaggregate error
rates by dialect and community, catching the over-flagging that silences minority
speakers~\cite{sap2019risk}.

\wordcount{47}

%%----------------------------------
\subsection{Cover Letter}
\label{subsec:Cover Letter}
%----------------------------------

A cover letter tailored to this advertisement is presented as a stand-alone letter in
Appendix~\ref{sec:Cover Letter}, formatted so that it could be sent unaltered with a job
application.

\wordcount{348 (see Appendix~\ref{sec:Cover Letter})}

%----------------------------------
\section{Data Set}
\label{sec:Data Set}
%----------------------------------

%%----------------------------------
\subsection{Description of the Data Sets}
\label{subsec:Description of the Data Sets}
%----------------------------------

Two public, account-free corpora of crowd-labelled user text underpin the analysis
(Table~\ref{tab:datasets}). Davidson et al.'s hate speech corpus~\cite{davidson2017automated}
holds 24,783 tweets, each judged by at least three CrowdFlower annotators and assigned one
of three classes: 5.77\% hate speech, 77.43\% offensive language, 16.80\% neither.
Wulczyn et al.'s Wikipedia Talk corpus~\cite{wulczyn2017ex} holds 159,686 discussion
comments carrying 1,598,289 toxicity judgements from 4,301 annotators, roughly ten per
comment; 9.62\% are toxic by majority vote. A companion file records each annotator's
demographics. Both download without an account from the URLs in Table~\ref{tab:datasets}.

\wordcount{90}

\begin{table}[h]
\centering
\caption{The two public corpora used in this case study. Class balance is computed from
the raw label files; the positive class is the item a moderator would be asked to act on.}
\label{tab:datasets}
\begin{tabular}{l>{\raggedright\arraybackslash}p{4.6cm}>{\raggedright\arraybackslash}p{4.6cm}}
\toprule
 & \textbf{Dataset A: Davidson et al.} & \textbf{Dataset B: Wulczyn et al.} \\
\midrule
Source & Twitter & English Wikipedia talk pages \\
URL & \url{https://github.com/t-davidson/hate-speech-and-offensive-language} &
      \url{https://figshare.com/articles/dataset/Wikipedia_Talk_Labels_Toxicity/4563973} \\
Instances & 24,783 tweets & 159,686 comments \\
File size & 2.4~MB (CSV) & 109~MB (two TSV files) \\
Attributes & \texttt{count}, \texttt{hate\_speech}, \texttt{offensive\_language},
             \texttt{neither}, \texttt{class}, \texttt{tweet} &
             \texttt{rev\_id}, \texttt{comment}, \texttt{year}, \texttt{logged\_in},
             \texttt{ns}, \texttt{sample}, \texttt{split} \\
Annotations & $\geq 3$ per tweet (median 3, max 9) & 1,598,289 judgements, 4,301
              annotators, $\approx$10 per comment \\
Label scheme & 3 classes: hate speech / offensive / neither & Binary toxicity per
               annotator, aggregated by majority \\
Positive class & Hate speech (5.77\%, $n=1{,}430$) & Toxic (9.62\%, $n=15{,}362$) \\
Mean length & 85.4 characters & 446.9 characters \\
Licence & Public research release (MIT) & CC0 1.0 Public Domain \\
\bottomrule
\end{tabular}
\end{table}

%%----------------------------------
\subsection{Rationale for Selection}
\label{subsec:Rationale for Selection}
%----------------------------------

Both corpora are exactly the modality the advertised role's classifiers consume: short,
user-generated, crowd-labelled text. Their label schemes deliberately disagree ---
Davidson separates hate speech from merely offensive language, while Wikipedia collapses
both into ``toxic'' --- which lets me test whether a triage model survives a change of
platform, the failure mode that matters in production. Limitations are real: both are
English-only, Davidson's lexicon-seeded sampling over-represents slurs, Wikipedia talk
pages are not video comments, neither includes audio or vision, and majority-vote labels
discard annotator disagreement.

\wordcount{83}

%----------------------------------
\section{Data Analysis}
\label{sec:Data Analysis}
%----------------------------------

\subsection{Method}
\label{subsec:Method}

Both corpora were reduced to the same operational question a moderation queue actually
asks: \emph{should this item be escalated to a human reviewer?} In the Davidson corpus the
positive class is \texttt{hate speech}; in the Wikipedia corpus it is a comment that the
majority of its annotators marked toxic. Text was represented as sublinear TF--IDF over
word unigrams and bigrams, capped at 20,000 features.

Two algorithms were compared: a \textbf{linear support vector
machine}~\cite{cortes1995support}, the standard strong baseline for high-dimensional
sparse text~\cite{joachims1998text}, and a \textbf{feed-forward neural network} with one
64-unit hidden layer, which can represent feature interactions a linear boundary cannot.
Neither was used in previous coursework. Both were implemented with
scikit-learn~\cite{pedregosa2011scikit} and evaluated under stratified five-fold cross
validation. Critically, the vectoriser sits inside the pipeline and is therefore refitted
on each training fold; fitting it once over the whole corpus would leak test-fold
vocabulary into training and inflate every metric reported
below~\cite{pageEvaluating}.

\subsection{Choice of Evaluation Metrics}
\label{subsec:Choice of Evaluation Metrics}

Accuracy is the wrong headline for this case study, and Table~\ref{tab:results} shows why
concretely. On the Davidson corpus the neural network is the \emph{more accurate} model
(94.5\% versus 92.2\%) while being the substantially worse moderation tool: it recovers
13.5\% of hate speech where the SVM recovers 41.1\%. With a 5.77\% positive class, a
model that predicted ``never escalate'' would score 94.2\% accuracy and be
useless~\cite{pageEvaluating}. Four metrics were therefore chosen against the
properties of this problem:

\begin{itemize}
    \item \textbf{Precision and recall} on the positive class, because the two error types
    have different costs and different owners. A false negative leaves harmful content in
    front of users; a false positive removes a legitimate creator's post and generates an
    appeal. Neither is captured by a single accuracy figure.

    \item \textbf{PR-AUC} as the primary threshold-independent score. Both corpora are
    dominated by negatives, the regime in which precision-recall curves are more
    informative than ROC curves, because ROC's false-positive rate is diluted by a very
    large negative denominator and so looks flattering~\cite{davis2006relationship,
    saito2015precision}.

    \item \textbf{ROC-AUC} reported alongside it, since it is insensitive to class
    distribution and therefore the fairer basis for comparing the same model across two
    corpora with different base rates~\cite{pageEvaluating}.

    \item \textbf{Precision@1\%}, the precision over the top 1\% of the ranked stream.
    This is the metric the deployed system is actually judged on: moderator headcount is
    fixed, so what matters is the purity of the queue a reviewer works through, not
    performance at an arbitrary 0.5 threshold.
\end{itemize}

Model comparison uses a paired two-tailed $t$-test over the per-fold differences rather
than a comparison of means, so that fold-to-fold variance is accounted
for~\cite{pageEvaluating}. The test is reported with the caveat that repeated-CV folds
are not fully independent, which makes it liberal~\cite{dietterich1998approximate}.

\begin{table}[h]
\centering
\caption{Five-fold cross-validated performance, averaged across folds. Bold marks the
better of the two models on each row. $p$-values are from paired two-tailed $t$-tests on
per-fold differences. Note that on Dataset~A the neural network is the more accurate model
yet materially the worse moderation tool.}
\label{tab:results}
\begin{tabular}{llcccc}
\toprule
& & \multicolumn{2}{c}{\textbf{Dataset A: Davidson}} & \multicolumn{2}{c}{\textbf{Dataset B: Wikipedia}} \\
\cmidrule(lr){3-4}\cmidrule(lr){5-6}
\textbf{Metric} & & \textbf{Linear SVM} & \textbf{Neural net} & \textbf{Linear SVM} & \textbf{Neural net} \\
\midrule
Accuracy      & & 0.922 & \textbf{0.945} & 0.938 & \textbf{0.957} \\
Precision     & & 0.353 & \textbf{0.619} & 0.640 & \textbf{0.851} \\
Recall        & & \textbf{0.411} & 0.135 & \textbf{0.820} & 0.674 \\
F1            & & \textbf{0.379} & 0.216 & 0.719 & \textbf{0.752} \\
ROC-AUC       & & \textbf{0.829} & 0.801 & 0.959 & 0.960 \\
PR-AUC        & & 0.337 & 0.327 & 0.839 & 0.840 \\
Precision@1\% & & 0.632 & 0.640 & \textbf{1.000} & 0.999 \\
\midrule
\multicolumn{6}{l}{\emph{Paired $t$-test, SVM $-$ neural net:}} \\
$p$ (F1)      & & \multicolumn{2}{c}{0.014$^{*}$} & \multicolumn{2}{c}{$<$0.001$^{*}$} \\
$p$ (PR-AUC)  & & \multicolumn{2}{c}{0.299} & \multicolumn{2}{c}{0.778} \\
$p$ (ROC-AUC) & & \multicolumn{2}{c}{0.012$^{*}$} & \multicolumn{2}{c}{0.651} \\
\bottomrule
\multicolumn{6}{l}{\footnotesize $^{*}$ significant at $\alpha = 0.05$. Base rates: 5.77\% (A), 9.62\% (B).}
\end{tabular}
\end{table}

\subsection{What the Algorithms Reveal}
\label{subsec:What the Algorithms Reveal}

\textbf{The two models rank almost identically; they differ only in where they cut.}
This is the least obvious and most useful result. On both corpora the difference in PR-AUC
is not statistically significant ($p = 0.299$ and $p = 0.778$), and on the Wikipedia
corpus neither is ROC-AUC ($p = 0.651$). Yet precision and recall differ significantly
every time. The two algorithms have learned an equivalent ordering of the data, and every
apparent difference in Table~\ref{tab:results} comes from the operating point rather than
from model capacity --- the SVM was given \texttt{class\_weight="balanced"} and so cuts
aggressively for recall, while the network, trained on the raw imbalance, cuts
conservatively for precision. Figure~\ref{fig:pr curves} shows this directly: the curves
sit on top of one another, and only the points chosen along them differ. For an engineer
this reframes the decision. The choice between these two models is not a modelling
question at all; it is a threshold policy set by how many moderators are rostered.

\textbf{The task is far harder on one corpus than the other.} On Wikipedia both models are
strong: PR-AUC $\approx 0.84$ against a 9.62\% base rate, and precision over the top 1\% of
the ranked stream is effectively perfect (1.000 and 0.999). A reviewer working the head of
that queue would see almost nothing but genuine toxicity. On Davidson both models are
weak: PR-AUC $\approx 0.33$, and precision@1\% of only 0.63. The gap is not a data-volume
artefact --- it is that separating \emph{hate speech} from \emph{merely offensive
language} is intrinsically harder than separating toxic from civil, which is precisely the
problem Davidson et al.\ set out to demonstrate~\cite{davidson2017automated}.

\begin{figure}[h]
    \centering
    \includegraphics[width=\linewidth]{img/fig_pr_curves.pdf}
    \Description{Two precision-recall panels. On the Davidson corpus both models trace
    nearly identical curves reaching PR-AUC of about 0.33 against a 5.8 percent base rate.
    On the Wikipedia corpus both trace nearly identical curves reaching PR-AUC of about
    0.84 against a 9.6 percent base rate. In both panels the two model curves overlap
    almost exactly.}
    \caption{Precision-recall curves, pooled across folds. The two algorithms are
    near-indistinguishable as \emph{rankers} on both corpora; the significant differences
    in Table~\ref{tab:results} come from threshold placement, not ranking quality.}
    \label{fig:pr curves}
\end{figure}

\begin{figure}[h]
    \centering
    \includegraphics[width=\linewidth]{img/fig_roc_curves.pdf}
    \Description{Two ROC panels. Both models reach ROC-AUC around 0.80 to 0.83 on the
    Davidson corpus and around 0.96 on the Wikipedia corpus, well above the random-guessing
    diagonal.}
    \caption{ROC curves, pooled across folds. ROC-AUC of 0.96 on Dataset~B looks
    excellent, but compare it with the PR-AUC of 0.84 in Figure~\ref{fig:pr curves}: with
    90\% of the corpus negative, the false-positive rate is diluted and ROC flatters the
    models~\cite{davis2006relationship}.}
    \label{fig:roc curves}
\end{figure}

\subsection{Are the Two Sources Complementary or Contradicting?}
\label{subsec:Are the Two Sources Complementary or Contradicting?}

\textbf{Both.} They agree about the models and disagree about the labels, and separating
those two claims is the substantive finding of this analysis.

They are \textbf{complementary about model selection}. Independently, on corpora that
differ by a factor of six in size, by platform, by mean document length and by annotation
protocol, both corpora return the same verdict: the two algorithms rank equivalently and
the SVM achieves it at roughly one-fifth of the training cost. A conclusion that
replicates across two unrelated sources is far stronger than one drawn from either alone.

They are \textbf{directly contradicting about what the label means}, and the contradiction
is severe enough to invalidate any attempt to pool them. Training on Wikipedia and testing
on Davidson, the SVM flags \textbf{73.3\% of the entire stream} against a true base rate of
5.8\%, and precision falls from 0.640 to 0.065 --- ten times worse than the within-corpus
figure. Figure~\ref{fig:transfer} shows why. The Wikipedia-trained SVM flags 82.7\% of
Davidson's hate speech and \emph{82.7\% of its merely offensive language} --- the identical
rate. The model has no discriminative power whatsoever between the two classes, because
the corpus it learned from never encoded a distinction between them. The neural network
behaves the same way, flagging 68.4\% of hate speech and 70.8\% of non-hate offensive
language, if anything inverting the ordering.

\begin{figure}[h]
    \centering
    \includegraphics[width=0.72\linewidth]{img/fig_transfer.pdf}
    \Description{A grouped bar chart. A Wikipedia-trained linear SVM flags 83 percent of
    Davidson hate speech and 83 percent of Davidson offensive-but-not-hate content, and 27
    percent of the neither category. The neural network flags 68 percent, 71 percent and 9
    percent respectively. Both models flag the hate and the merely-offensive categories at
    essentially the same rate.}
    \caption{Applying a Wikipedia-trained model to Davidson's data. Both models flag
    genuine hate speech and merely offensive language at effectively the same rate ---
    they cannot tell the two apart, because the Wikipedia label scheme never separated
    them. This is invisible to any single-corpus evaluation.}
    \label{fig:transfer}
\end{figure}

The reverse direction fails differently and just as informatively: a Davidson-trained
neural network flags only 0.8\% of Wikipedia against a 9.6\% base rate, recovering 5.4\% of
genuine toxicity. Trained on a narrow, lexicon-seeded definition, it does not recognise
most of what Wikipedia's annotators consider harmful.

\textbf{Do all data sources give insights? Not equally.} Wikipedia supports a deployable
triage model. Davidson does not --- but its value is diagnostic rather than predictive: it
is the only one of the two that can expose the over-flagging failure at all, precisely
because it preserves the distinction the other corpus collapses. Discarding it as the
``worse'' dataset on PR-AUC would have discarded the finding.

\subsection{Limitations and Implications for Practice}
\label{subsec:Limitations and Implications for Practice}

Three limitations bound these conclusions. Both corpora are English-only and text-only,
whereas the advertised role concerns short-form video across many languages. Majority-vote
aggregation discards annotator disagreement, and disagreement is signal --- an item ten
annotators split evenly on is a different object from one they unanimously cleared. And
the paired $t$-test over five overlapping folds is liberal, so the significant results
above should be read as indicative rather than conclusive~\cite{dietterich1998approximate}.

A further concern is one the metrics here cannot see. Published work shows that classifiers
trained on corpora of this kind systematically over-flag African-American
English~\cite{sap2019risk, blodgett2016demographic}, and the aggregate PR-AUC that looks
healthy on Dataset~B can conceal a much worse error rate for a specific community. The
Wikipedia corpus ships annotator demographics for exactly this reason. Before deployment,
the evaluation in Table~\ref{tab:results} would need to be disaggregated by dialect and
speaker group, not merely repeated at higher precision.

The practical recommendation follows from the first result. Ship the linear SVM: it ranks
as well as the neural network on every threshold-independent measure, trains in roughly a
fifth of the time, and its coefficients can be inspected when a creator appeals a removal.
Then treat the decision threshold as an operational control tied to review capacity rather
than a fixed model property. Most importantly, never validate such a model on a single
labelled source. A system reporting 0.96 ROC-AUC in-corpus flagged three-quarters of a
different platform's traffic, and no amount of cross validation \emph{within} one dataset
would have revealed it~\cite{vidgen2020directions}.


\clearpage

\begin{center}
    \noindent \textcolor{purple}{\emph{-- Content below this line does not count for the page/word limit. --}}
\end{center}

\section*{Generative AI Attribution and Declaration}

\noindent\textbf{AI Attribution Statement} (generated with the AI Attribution Toolkit,
\url{https://aiattribution.github.io/}):

\begin{quote}
\emph{This work was produced with the assistance of generative AI. AI was used to draft
and edit prose, to write and debug the analysis code, and to structure the document.
Dataset selection, the framing of the case study, the interpretation of results, and all
final judgements are the author's own. All AI-assisted output was reviewed, tested against
the data, and corrected by the author.}
\end{quote}

\noindent Specifically, Claude (Anthropic) was used to: locate and verify the public
corpora and their access URLs; draft the Python analysis pipeline in
Appendix~\ref{sec:Link to GitHub Repository}; draft and condense prose against the word
limits in this specification; and format the \LaTeX{} and \BibTeX{} sources. Every numeric
result reported in Section~\ref{sec:Data Analysis} was produced by executing the code in
the linked repository, not by the language model. A record of the prompts, intermediate
outputs and issues encountered is given in Appendix~\ref{sec:Development Process}.

\begin{acks}
I acknowledge the people of the Woi Wurrung and Boon Wurrung language groups of the
eastern Kulin Nation, on whose unceded lands RMIT University's Melbourne campuses stand. I
respectfully recognise their Ancestors and Elders, past and present.

I thank Thomas Davidson and colleagues, and Ellery Wulczyn and colleagues, for releasing
the corpora that made this analysis possible.
\end{acks}

\bibliographystyle{ACM-Reference-Format}
\bibliography{99-references}

\appendix
%----------------------------------
\section{Job Listing}
\label{sec:Job Listing}
%----------------------------------

The advertisement analysed in Section~\ref{sec:Data Science Role} is reproduced below as
retrieved on 14 August 2026 from
\url{https://au.gradconnection.com/employers/tiktok/jobs/tiktok-machine-learning-engineer-graduate-trust-and-safety-engineering-2027-start-5/}.
Applications close 30 August 2026.

\begin{figure}[h]
    \centering
    \fbox{\includegraphics[width=0.92\linewidth, page=1]{img/job_advertisement.pdf}}
    \Description{A one-page copy of the TikTok job advertisement for a Machine Learning
    Engineer Graduate in Trust and Safety Engineering, listing the team introduction,
    nine responsibilities, four minimum qualifications and two preferred qualifications.}
    \caption{Job advertisement: TikTok, Machine Learning Engineer Graduate (Trust and
    Safety Engineering) --- 2027 Start, Sydney NSW.}
    \label{fig:job advertisement}
\end{figure}

\clearpage

%----------------------------------
\section{Development Process}
\label{sec:Development Process}
%----------------------------------

This appendix documents how the executive summary was produced, as evidence of process
rather than of outcome alone.

\subsection{Pipeline}

The analysis is five scripts, run in order, each writing artefacts the next consumes
(Table~\ref{tab:pipeline}). Every number quoted in Section~\ref{sec:Data Analysis} comes
from \texttt{results/results.json}, which is regenerated from scratch by
\texttt{03\_analysis.py}.

\begin{table}[h]
\centering
\caption{Analysis pipeline. Total runtime on a consumer laptop was approximately
25 minutes, dominated by the neural network folds on the larger corpus.}
\label{tab:pipeline}
\begin{tabular}{llp{6.2cm}}
\toprule
\textbf{Step} & \textbf{Script} & \textbf{Purpose} \\
\midrule
1 & \texttt{01\_download\_data.py} & Fetch both corpora from GitHub and Figshare; no
                                     account required \\
2 & \texttt{02\_profile\_data.py}  & Record exact row counts, attributes and class balance
                                     quoted in Section~\ref{sec:Data Set} \\
3 & \texttt{03\_analysis.py}       & Stratified 5-fold cross validation of both models on
                                     both corpora, cross-corpus transfer, paired $t$-tests \\
4 & \texttt{04\_job\_ad\_pdf.py}   & Typeset the retrieved advertisement as
                                     Figure~\ref{fig:job advertisement} \\
5 & \texttt{05\_figures.py}        & Render the precision-recall and ROC curves and the
                                     transfer chart \\
\bottomrule
\end{tabular}
\end{table}

\subsection{Core Modelling Code}

Listing~\ref{lst:models} shows the two estimators and the feature pipeline. The
vectoriser sits \emph{inside} the pipeline so that
\texttt{StratifiedKFold} refits it on each training fold; fitting it once over the whole
corpus would leak test-fold vocabulary and inflate every metric, the pitfall
\citet{pageEvaluating} raises on slide~35.

\begin{lstlisting}[caption={Feature pipeline and the two estimators
(\texttt{code/03\_analysis.py}).}, label={lst:models}]
def make_vectoriser():
    # Fitted inside every fold, never on the full corpus.
    return TfidfVectorizer(
        sublinear_tf=True, min_df=3, max_features=20000,
        ngram_range=(1, 2), strip_accents="unicode", lowercase=True,
    )

def make_models():
    return {
        "Linear SVM": Pipeline([
            ("tfidf", make_vectoriser()),
            ("clf", LinearSVC(C=1.0, class_weight="balanced",
                              max_iter=5000, random_state=SEED)),
        ]),
        "Neural network (MLP)": Pipeline([
            ("tfidf", make_vectoriser()),
            ("clf", MLPClassifier(hidden_layer_sizes=(64,), activation="relu",
                                  solver="adam", alpha=1e-4, batch_size=256,
                                  learning_rate_init=1e-3, max_iter=30,
                                  early_stopping=True, n_iter_no_change=3,
                                  validation_fraction=0.1, random_state=SEED)),
        ]),
    }

def precision_at_capacity(y_true, score, capacity=0.01):
    """Precision over the top-scoring slice a moderation queue can review."""
    k = max(1, int(round(len(y_true) * capacity)))
    top = np.argsort(score)[::-1][:k]
    return float(y_true[top].mean()), k
\end{lstlisting}

\subsection{Issues Encountered and How They Were Addressed}

\begin{enumerate}
    \item \textbf{The first advertisement selected had closed.} The 2026-start version of
    this role expired before submission. Rather than cite a dead link, the employer's job
    index was searched for the live equivalent; the 2027-start listing was confirmed open
    (an active \emph{Apply} control, closing 30 August 2026) and substituted. The text of
    the two postings is materially identical.

    \item \textbf{The obvious toxicity dataset was behind a login.} The widely used Jigsaw
    release is distributed through Kaggle, which requires an account and so fails the
    specification's requirement that the link be accessible. The original Wulczyn et al.
    release on Figshare was used instead: same underlying annotations, CC0 licence, no
    account.

    \item \textbf{Encoded whitespace polluted the features.} Wikipedia comments store line
    and tab breaks as the literal strings \texttt{NEWLINE\_TOKEN} and \texttt{TAB\_TOKEN}.
    Left in place these become two of the highest-frequency terms in the vocabulary and
    carry no signal. They are stripped before vectorising.

    \item \textbf{The two estimators expose different score interfaces.}
    \texttt{LinearSVC} provides \texttt{decision\_function} but not
    \texttt{predict\_proba}. Because both ROC and precision-recall curves depend only on
    the \emph{ranking} induced by the score, the raw decision values are used directly;
    no calibration step is needed for the curves, though it would be needed before
    interpreting a score as a probability in production.

    \item \textbf{Neural network runtime.} A dense hidden layer over a full-vocabulary
    sparse matrix was too slow to cross-validate on the larger corpus. The vocabulary was
    capped at 20,000 features and the network reduced to a single 64-unit hidden layer
    with early stopping. Both models share the identical vectoriser configuration so the
    comparison stays fair.

    \item \textbf{Class imbalance made the default threshold useless for the SVM.} At the
    default decision boundary the minority class was barely predicted at all.
    \texttt{class\_weight="balanced"} was enabled, which is why the SVM's operating point
    favours recall over precision --- a choice revisited in
    Section~\ref{sec:Data Analysis}.
\end{enumerate}

\subsection{Use of Generative AI}

Claude (Anthropic) was used throughout, in the manner declared on
page~\pageref{sec:Data Analysis}. Representative prompts:

\begin{itemize}
    \item ``Find live machine learning engineer job advertisements in Australia and verify
    the listing URLs actually resolve.''
    \item ``Propose two publicly available datasets relevant to content-moderation triage
    that download without an account, and confirm the download URLs respond.''
    \item ``Write a stratified cross-validation harness comparing a linear SVM and an MLP,
    fitting the vectoriser inside each fold, and report precision, recall, F1, ROC-AUC and
    PR-AUC plus a paired $t$-test across folds.''
    \item ``Condense this paragraph to at most 50 words without losing the citation.''
\end{itemize}

Two defects survived into a compiled draft and were caught only by rendering the PDF and
inspecting every page rather than trusting the source. The typeset job advertisement
originally ran to two pages, of which Figure~\ref{fig:job advertisement} included only the
first --- the preferred qualifications were silently missing. Separately, the cover letter
inherited the ACM template's one-and-a-half line spacing and broke across two pages,
defeating the purpose of a stand-alone letter. Both are corrected in this version. The
lesson generalises to the analysis itself: generated output looks finished long before it
is correct, and the only reliable check is to inspect the artefact rather than the code
that produced it.

%----------------------------------
\section{Link to GitHub Repository}
\label{sec:Link to GitHub Repository}
%----------------------------------

All code, the downloaded-data manifest and the generated results are available at:

\url{https://github.com/heetchanchad15/Case-Studies-Assignment}

\clearpage

%----------------------------------
\section{Cover Letter}
\label{sec:Cover Letter}
%----------------------------------

The letter below is written to stand alone: it carries its own sender and recipient
blocks and no report formatting, so it could be sent unaltered with an application.
\wordcount{348}

\newpage

\thispagestyle{empty}
\begingroup
\singlespacing % a business letter is single-spaced, unlike the ACM manuscript body

\noindent
{\large\textbf{Heet Chanchad}}\\[2pt]
Melbourne VIC, Australia\\
s4218211@student.rmit.edu.au

\vspace{14pt}

\noindent 14 August 2026

\vspace{12pt}

\noindent
Early Careers Team\\
Trust \& Safety Engineering\\
TikTok Australia\\
Sydney NSW

\vspace{12pt}

\noindent\textbf{Re: Machine Learning Engineer Graduate (Trust and Safety Engineering) ---
2027 Start}

\vspace{10pt}

\noindent Dear Hiring Team,

\vspace{8pt}

\noindent I am writing to apply for the Machine Learning Engineer Graduate role in Trust
and Safety Engineering. I am completing a Master of Data Science at RMIT University's City
campus, which I began in March 2026 and will finish in March 2028, and I am available to
begin from that date. I want to spend my career building systems that reduce harm at
scale, not systems that maximise attention, and your team is one of the few places where
that is the explicit brief.

\vspace{8pt}

\noindent The advertisement asks for someone who can ``develop and build up
highly-scalable classifiers'' and ``work on model iterations''. My most recent project is
that problem in miniature. I built a content-triage pipeline over two public corpora ---
24,783 crowd-labelled tweets and 159,686 Wikipedia Talk comments carrying 1.6 million
toxicity judgements --- and compared a linear support vector machine against a
feed-forward neural network under stratified five-fold cross validation.

\vspace{8pt}

\noindent Two findings from that work speak directly to your team's problem. First, I
evaluated the models the way a moderation queue actually runs. Because only 5.8\% and
9.6\% of the two corpora are removal-worthy, accuracy is close to meaningless, so I
reported precision, recall and area under the precision-recall curve, and measured
precision across the top 1\% of the ranked stream --- the slice a fixed review team could
genuinely process in a day. Second, I tested whether a model survives a change of source.
A classifier scoring 0.96 ROC-AUC on its own corpus flagged 73\% of the other platform's
traffic against a true base rate of 5.8\%, because the two label schemes disagree about
whether offensive language is the same thing as hate speech. That failure is invisible in
single-dataset metrics and would surface in production as mass wrongful removals of
ordinary speech.

\vspace{8pt}

\noindent I also treat fairness as part of the engineering rather than a review step after
it. Published work shows hate-speech classifiers systematically over-flag African-American
English, and the Wikipedia corpus ships annotator demographics precisely because who
labels the data decides what counts as toxic. Before I would ship a classifier, I would
want its error rates broken down by dialect and community, not just its headline score.

\vspace{8pt}

\noindent I would welcome the opportunity to discuss how I could contribute to the team.
Thank you for your consideration.

\vspace{12pt}

\noindent Yours sincerely,

\vspace{16pt}

\noindent Heet Chanchad

\endgroup


\end{document}
